\section{限界}
\label{sec:limitations}
命題\ref{prop:noncausality}はモデル内の入力非依存性を保証する一方、真の構造の回復は変数集合、ラグ数、表現力、標本数に依存し、相関する入力が予測を代替すれば選択辺も変わり得る。Lorenz96で確認したのは、独立した予備系列の正解構造で設定を選んだ条件での性能である。未知構造で同等の回復を実現できるとは限らず、予測損失の小ささだけでは選択グラフの正しさを確認できない。

Favoritaでは単一店舗の12カテゴリと一つの学習初期値に対象が限られる。販促、価格、祝日、在庫、来店者数などを明示的に条件付けていないため、選択辺が共通要因を反映する可能性がある。曜日調整後も残る季節性や非定常性も推定に影響する。初期値や店舗を変えた再現性、外部情報を加えた構造の安定性は、本稿の結果からは判断できない。

有効係数とヤコビアンの対応は観測状態上で調べられ、介入効果の同定には追加の因果仮定を要する。対象・ラグ別の生成器が表す交互作用の範囲は同一ラグ内に限られる。
